\section*{Capítulo \ref{ch:g9:arith} --- Aritmética: divisores y números primos}

\begin{solution}{exo:g9:arith:1}
\omterm{def:g9:arith:divisor}{Divisores} de $36$: $1, 2, 3, 4, 6, 9, 12, 18, 36$.
\omterm{def:g9:arith:divisor}{Divisores} de $45$: $1, 3, 5, 9, 15, 45$.
\omterm{def:g9:arith:divisor}{Divisores} de $17$: solo $1$ y $17$ ($17$ es primo).
\end{solution}

\begin{solution}{exo:g9:arith:2}
$2\,346$ termina en $6$: \omterm{def:g6:wholes:divisible}{divisible} entre $2$, no entre $5$. \omterm{def:g1:addition:def}{Suma} de cifras
$2 + 3 + 4 + 6 = 15$: \omterm{def:g6:wholes:divisible}{divisible} entre $3$, no entre $9$. Las dos últimas cifras son
$46$, y $46 = 4 \times 11 + 2$ no es \omterm{def:g6:wholes:divisible}{divisible} entre $4$: $2\,346$ no es
\omterm{def:g6:wholes:divisible}{divisible} entre $4$.
\end{solution}

\begin{solution}{exo:g9:arith:3}
$72 = 2^3 \times 3^2$; \quad
$150 = 2 \times 3 \times 5^2$; \quad
$210 = 2 \times 3 \times 5 \times 7$; \quad
$121 = 11^2$.
\end{solution}

\begin{solution}{exo:g9:arith:4}
$101$: prueba con los primos $p$ tales que $p^2 \leq 101$, es decir $2, 3, 5, 7$. Ninguno
\omterm{def:g9:arith:divisor}{divide} a $101$ (es \omterm{def:g3:numbers:evenodd}{impar}, \omterm{def:g1:addition:def}{suma} de cifras $2$, no termina en $0/5$,
$101 = 7 \times 14 + 3$): $101$ es primo.

$91 = 7 \times 13$: no es primo.

$143 = 11 \times 13$: no es primo.
\end{solution}

\begin{solution}{exo:g9:arith:5}
\omterm{def:g9:arith:divisor}{Divisores} comunes de $48$ y $60$: los \omterm{def:g9:arith:divisor}{divisores} de 48 son $1, 2, 3, 4, 6, 8,
12, 16, 24, 48$; los \omterm{def:g9:arith:divisor}{divisores} de $60$ son $1, 2, 3, 4, 5, 6, 10, 12, 15, 20,
30, 60$; los comunes son $1, 2, 3, 4, 6, 12$, así que $\gcd(48,60) = 12$.

Por descomposición: $48 = 2^4 \times 3$ y $60 = 2^2 \times 3 \times 5$;
primos comunes con los \omterm{def:g8:powers:def}{exponentes} menores: $2^2 \times 3 = 12$.
\end{solution}

\begin{solution}{exo:g9:arith:6}
$\gcd(255, 154)$:
\begin{align*}
255 &= 154 \times 1 + 101, \\
154 &= 101 \times 1 + 53, \\
101 &= 53 \times 1 + 48, \\
53 &= 48 \times 1 + 5, \\
48 &= 5 \times 9 + 3, \\
5 &= 3 \times 1 + 2, \\
3 &= 2 \times 1 + 1, \\
2 &= 1 \times 2 + 0 .
\end{align*}
Último \omterm{def:g3:division:remainder}{resto} no nulo: $\gcd(255, 154) = 1$ (son \omterm{def:g9:arith:gcd}{primos entre sí}).

$\gcd(1053, 325)$:
\begin{align*}
1053 &= 325 \times 3 + 78, \\
325 &= 78 \times 4 + 13, \\
78 &= 13 \times 6 + 0 .
\end{align*}
$\gcd(1053, 325) = 13$.
\end{solution}

\begin{solution}{exo:g9:arith:7}
Algoritmo de Euclides: $588 = 504 \times 1 + 84$; $504 = 84 \times 6 + 0$:
$\gcd(588, 504) = 84$. Entonces
\[
\frac{588}{504} = \frac{588 \div 84}{504 \div 84} = \frac{7}{6},
\]
que es irreducible.
\end{solution}

\begin{solution}{exo:g9:arith:8}
El número de ramos tiene que dividir a $84$ y a $126$; el mayor
posible es $\gcd(84, 126)$. Descomposiciones: $84 = 2^2 \times 3 \times
7$, $126 = 2 \times 3^2 \times 7$, así que el MCD es $2 \times 3 \times 7 =
42$. Puede hacer $42$ ramos, cada uno con
$\frac{84}{42} = 2$ rosas y $\frac{126}{42} = 3$ tulipanes.
\end{solution}

\begin{solution}{exo:g9:arith:9}
Hace falta el mínimo común \omterm{def:g9:arith:divisor}{múltiplo}. $24 = 2^3 \times 3$ y
$36 = 2^2 \times 3^2$; tomando cada primo con el \emph{mayor}
\omterm{def:g8:powers:def}{exponente}: $\lcm = 2^3 \times 3^2 = 72$. Los transbordadores vuelven a salir juntos
$72$ minutos después de las 8:00, a las 9:12.
\end{solution}

\begin{solution}{exo:g9:arith:10}
\emph{1.} Todo \omterm{def:g9:arith:divisor}{divisor} común $d$ de $n$ y $n+1$ \omterm{def:g9:arith:divisor}{divide} también a su
\omterm{ex:g1:subtraction:difference}{diferencia} $(n+1) - n = 1$, luego $d = 1$: $\gcd(n, n+1) = 1$.

\emph{2.} Una \omterm{def:g9:fractions:fraction}{fracción} es irreducible exactamente cuando su \omterm{def:g4:fractions:def}{numerador} y su
\omterm{def:g4:fractions:def}{denominador} son \omterm{def:g9:arith:gcd}{primos entre sí}, que es el caso de $n$ y $n + 1$ por el apartado
1.
\end{solution}

\begin{solution}{pb:g9:arith:1}
\textbf{1.} Llena la de $3$ y viértela en la de $5$ (contenidos
$0/3 \to 3$ en la grande). Vuelve a llenar la de $3$ y viértela en la
de $5$ hasta que esté llena: la jarra grande solo admite $2$ más, y quedan
\[
3 - 2 = 1 \text{ L en la jarra pequeña.}
\]
Movimientos: llena $3$; vierte $3 \to 5$; llena $3$; vierte $3 \to 5$.

\textbf{2.} Llena la de $5$; viértela en la de $3$ (quedan $2$ en la
grande); vacía la de $3$; vierte los $2$ en la de $3$; llena la de $5$;
viértela en la de $3$ hasta llenarla --- admite $1$, y quedan
$\mathbf{4}$ L en la jarra grande. Seis movimientos.

\textbf{3.} Todas: $1$ (pregunta 1), $2$ (tras dos movimientos
de la pregunta 2), $3$ y $5$ (con un solo llenado), $4$ (pregunta
2), $6 = 3 + 3$ (la jarra pequeña llena más $3$ vertidos en la
grande), $7 = 5 + 2$, $8 = 5 + 3$ (las dos llenas). Toda cantidad entera
de $1$ a $8$ L se puede medir con la de $5$ y la de $3$.

\textbf{4.} Al principio las dos jarras contienen $0$, \omterm{def:g9:arith:divisor}{múltiplo} de $2$.
Llenar pone un contenido a $6$ o a $4$: \omterm{def:g3:numbers:evenodd}{par}. Vaciar lo pone a
$0$: \omterm{def:g3:numbers:evenodd}{par}. Verter mueve algo de agua entre jarras cuyos contenidos
eran \omterm{def:g3:numbers:evenodd}{pares}, y la cantidad vertida es \omterm{ex:g1:subtraction:difference}{diferencia} de \omterm{def:g3:numbers:evenodd}{números
pares} (el hueco libre o la cantidad disponible): todos los contenidos siguen siendo
\omterm{def:g3:numbers:evenodd}{pares} para siempre. Un objetivo \omterm{def:g3:numbers:evenodd}{impar} como $1$ L es inalcanzable.

\textbf{5.} $\gcd(6, 4) = 2$: solo cantidades \omterm{def:g3:numbers:evenodd}{pares} --- lo confirma la
pregunta 4. $\gcd(5, 3) = 1$: la ley permite toda cantidad entera,
y la pregunta 3 las realizó todas. El MCD es exactamente
la unidad de medida de las jarras.

\textbf{6.} $91 = 1 \times 65 + 26$; $65 = 2 \times 26 + 13$;
$26 = 2 \times 13 + 0$: $\gcd(91, 65) = 13$. Y
$2\,026 = 44 \times 46 + 2$; $46 = 23 \times 2 + 0$:
$\gcd(2\,026, 46) = 2$.

\textbf{7.} Vertiendo la de $5$ en la de $13$ una y otra vez: tras dos
llenados la jarra grande contiene $10$; en el tercer llenado solo caben
$3$, y quedan $5 - 3 = 2$ en la jarra pequeña --- el \omterm{def:g3:division:remainder}{resto} de
$13$ entre $5$ era $3$, y las cantidades $3$ (el hueco) y $2$
(lo sobrante) son exactamente los números de Euclides ($13 = 2 \times 5 + 3$,
$5 = 1 \times 3 + 2$). Siguiendo, aparece $3 - 2 = 1$: el
\omterm{def:g3:division:remainder}{resto} siguiente del algoritmo. La fuente ejecuta con agua las
divisiones de Euclides.

\textbf{8.} $\gcd(13, 5) = 1$, así que la ley de la pregunta 5 permite
toda cantidad entera --- y la cascada de la pregunta 7 produjo
de hecho $1$ L. Sí.

\textbf{9.} Un \omterm{def:g9:arith:divisor}{divisor} común de $n$ y $2n + 1$ \omterm{def:g9:arith:divisor}{divide} a
$2n + 1 - 2 \times n = 1$: tiene que ser $1$. Por tanto
$\gcd(n, 2n+1) = 1$ siempre.

\textbf{10.} Autobús A: 7:12, 7:24, 7:36, 7:48, 8:00\,\dots; autobús B:
7:18, 7:36, 7:54\,\dots\ Primera salida común: 7:36, al cabo de
$36$ minutos --- el primer \omterm{def:g9:arith:divisor}{múltiplo} común de $12$ y $18$.
Ley: $36 \times \gcd(12, 18) = 36 \times 6 = 216 = 12 \times
18$. Para $5$ y $3$: primer \omterm{def:g9:arith:divisor}{múltiplo} común $15$, y
$15 \times \gcd(5,3) = 15 \times 1 = 15 = 5 \times 3$.

\textbf{11.} Con un ciclo de $17$ años: la próxima coincidencia es
el primer \omterm{def:g9:arith:divisor}{múltiplo} común de $17$ y $4$; como
$\gcd(17, 4) = 1$, eso son $17 \times 4 = 68$ años --- las
cigarras se topan con el máximo una vez de cada cuatro emergencias. Con un
ciclo de $16$ años: $16$ es \omterm{def:g9:arith:divisor}{múltiplo} de $4$, así que \emph{cada}
emergencia cae en un máximo. Una duración de ciclo prima no comparte ningún factor
con ningún ciclo de depredador más corto, lo que separa las coincidencias todo lo
posible: la aritmética como camuflaje.

\textbf{12.} Cuatro opciones de \omterm{def:g8:powers:def}{exponente} para $2$, tres para $3$
y dos para $5$: $4 \times 3 \times 2 = 24$ \omterm{def:g9:arith:divisor}{divisores}.

\textbf{13.} Si $m = 2^{a} \times 3^{b} \times \cdots$, entonces
$m^2 = 2^{2a} \times 3^{2b} \times \cdots$: todos los \omterm{def:g8:powers:def}{exponentes} quedan
duplicados, luego \omterm{def:g3:numbers:evenodd}{pares}. En $360 = 2^3 \times 3^2 \times 5$, los
\omterm{def:g8:powers:def}{exponentes} de $2$ y de $5$ son \omterm{def:g3:numbers:evenodd}{impares}: $360$ no es un cuadrado
perfecto.

\textbf{14.} Un \omterm{def:g9:arith:divisor}{divisor} es su propio compañero exactamente cuando
$d = \frac nd$, es decir $n = d^2$: solo los cuadrados tienen un
\omterm{def:g9:arith:divisor}{divisor} central así. Para todos los demás $n$, los \omterm{def:g9:arith:divisor}{divisores} se reparten en
parejas, una cantidad \omterm{def:g3:numbers:evenodd}{par}. Así que: \omterm{def:g3:numbers:evenodd}{número impar} de \omterm{def:g9:arith:divisor}{divisores}
$\Leftrightarrow$ cuadrado perfecto. Comprobación: $36$ tiene por \omterm{def:g9:arith:divisor}{divisores}
$1, 2, 3, 4, 6, 9, 12, 18, 36$ --- nueve, \omterm{def:g3:numbers:evenodd}{impar}, y
$36 = 6^2$; mientras que $360$ tiene $24$ (pregunta 12), \omterm{def:g3:numbers:evenodd}{par}, y no es
cuadrado (pregunta 13).

\textbf{15.} A la taquilla $n$ le cambia el estado una vez cada alumno $k$
cuyo número \omterm{def:g9:arith:divisor}{divide} a $n$: en total, tantas veces como \omterm{def:g9:arith:divisor}{divisores}
tenga $n$. Una taquilla acaba \emph{abierta} cuando se le cambia el estado un \omterm{def:g3:numbers:evenodd}{número
impar} de veces --- por la pregunta 14, exactamente cuando $n$ es un
cuadrado perfecto. Taquillas abiertas: $1, 4, 9, 16, 25, 36, 49, 64,
81, 100$ --- diez.
\end{solution}
